The Brazilian semiarid region is characterized by high rainfall variability, significant evaporation losses and recurrent droughts, which make reservoir operation relevant to water security. This study aimed to develop and evaluate a web-based decision support system built on monthly water-balance simulation for analysing the operation of single reservoirs and multi-reservoir systems in Ceará, Brazil. The simulator accounts for inflow, evaporation computed from the mean water-surface area, withdrawal, spill and transfers; represents individual, series and parallel operation; and applies monthly rule curves with rationing. The interface, developed in React and connected to a FastAPI service and an SQLite database with 155 reservoirs, allows users to configure analyses, inspect indicators and monthly series, and export results. The implementation was verified against an independent spreadsheet calculation, through checks of the rationing rules, and through nine controlled test cases, covering 25 reservoir systems and 33,108 monthly balances. The spreadsheet and the simulator agreed for 24 systems, with maximum differences on the order of 10\textsuperscript{-14}~hm\textsuperscript{3}; the only discrepancy, at Canoas, was caused by a different area-volume curve in the spreadsheet and disappeared when the system database curve was used. All controlled cases passed. In the sensitivity analysis, demand was the most influential parameter for supply failures, evaporation had a moderate effect, and the initial storage had a negligible effect over the long record; the transfer trigger and transfer rate showed safety thresholds for the receiving reservoir. The applications showed that the rule curves eliminated the 142 failures that would occur at Mundaú without rationing, at the cost of 332 months of reduced withdrawal; that the failures in the Carnaubal--Batalhão system occurred when the secondary reservoir lacked storage to take over the joint demand; and that, in the Fogareiro--Quixeramobim system, the trigger rule caused transfers to switch on and off in consecutive months. The tests showed consistency between the implemented equations and the outputs produced in the controlled cases; since no comparison with observed storage was performed, the results do not constitute validation of how well the model represents actual operation.

% Separe as Keywords por ponto e vírgula.
\keywords{water resources; reservoir operation; decision support system; hydrological simulation; model verification.}
